% HUONG DAN SU DUNG SACH
\chapter*{Hướng Dẫn Sử Dụng Sách}
\addcontentsline{toc}{chapter}{Hướng Dẫn Sử Dụng Sách}
\markboth{Hướng Dẫn Sử Dụng Sách}{Hướng Dẫn Sử Dụng Sách}

\begin{truyen}{Quyển Sách Này Dành Cho Ai?}{Who Is This Book For?}
\chuhoa{Q}{uyển sách này phù hợp cho ba nhóm người đọc.}
\textit{(This volume is suitable for three groups of readers.)}

\textbf{Cha mẹ:} để soi lại cách yêu thương, nhắc nhở và kỳ vọng của mình với con.
\textit{(\textbf{Parents:} to reflect on how they love, guide, and place expectations on their children.)}

\textbf{Học sinh:} để hiểu phía sau nhiều lời la rầy đôi khi là nỗi lo chưa biết nói cho đúng.
\textit{(\textbf{Students:} to see that behind many scoldings there is often worry that does not know how to speak well.)}

\textbf{Giáo viên:} để có thêm chất liệu kết nối giữa nhà trường và gia đình trong các buổi sinh hoạt hoặc tư vấn.
\textit{(\textbf{Teachers:} to gain material for building connection between school and family in class meetings or counselling.)}
\end{truyen}

\vspace{6pt}

\begin{tcolorbox}[
  colback=clayrose!8,
  colframe=clayrose!70!black,
  coltitle=black,
  fonttitle=\bfseries,
  title={Giải Thích Các Ô Trong Sách \quad\textit{\small(Book Guide)}},
  arc=4pt, boxrule=1pt, breakable, enhanced
]

\smallskip
\textbf{Ô Xanh Đậm --- Màn Đối Thoại Chính (\textit{Main Dialogue})}\\
Đây là các cuộc nói chuyện rất gần với đời sống trong nhà. Đọc chậm để thấy đâu là lời thương, đâu là lời làm đau, và đâu là phần hai bên chưa nghe nhau đủ.\\
\textit{(These are dialogues close to real family life. Read slowly to notice where love appears, where pain is caused, and where both sides still fail to hear each other fully.)}

\medskip
\textbf{Ô Đen Xám --- Điều Còn Lại Trong Nhà (\textit{Takeaway})}\\
Một hoặc hai câu ngắn giữ lại tinh thần cốt lõi của đoạn đối thoại.\\
\textit{(One or two short sentences preserving the core of the dialogue.)}

\medskip
\textbf{Ô Xanh Nhạt --- Easy English}\\
Hai câu tiếng Anh ngắn để nhớ ý chính và luyện thêm cách diễn đạt.\\
\textit{(Two short English lines to remember the key idea and practise expression.)}

\medskip
\textbf{Ô Vàng --- Gợi Ý Cùng Nghĩ (\textit{Shared Reflection})}\\
Hai câu hỏi ngắn để cha mẹ và con có thể cùng nghĩ hoặc cùng nói với nhau sau khi đọc.\\
\textit{(Two short questions for parents and children to reflect on or discuss together after reading.)}

\smallskip
\end{tcolorbox}

\vspace{6pt}

\begin{baihoc}
Gợi ý cách đọc hiệu quả nhất: mỗi ngày một vài đoạn, tốt nhất đọc chậm và dừng lại ở đoạn làm mình thấy chạm. Đọc quyển này không phải để thắng nhau, mà để nói với nhau tốt hơn.
\textit{(Best reading method: a few pieces each day, slowly, stopping at the ones that touch you. This book is not for winning against each other, but for speaking better with each other.)}
\end{baihoc}

\begin{ghinhoanh}
\textbf{Reading tips:}

Read slowly. Some short dialogues stay long.

Do not read to win. Read to understand.
\end{ghinhoanh}

\ngancach